\subsection{问题三：预报更新下的购电修订与滚动调度}
\subsubsection{调整机会与结算模型}
新预报改变了未来供需的判断，也使已提交的计划可能偏多或偏少。问题三允许在6:00、12:00和18:00修订当天尚未交付的购电量，但减少和增加购电都需要按约定结算。因此，每次更新都应比较维持计划与调整计划的费用。

设$g_t$为0点原计划，$a_t$为最终有效购电量，定义
\begin{equation}
\delta_t^+=(a_t-g_t)^+,\qquad \delta_t^-=(g_t-a_t)^+.
\end{equation}
取消部分退回原购电费并支付50\%违约费，新增部分按1.5倍电价购买。最终实际费用为
\begin{equation}
C_{3,d}=\sum_t\left[p_tg_t-0.5p_t\delta_t^-+1.5p_t\delta_t^++5p_te_t\right].
\label{eq:q3bill}
\end{equation}
例如电价为1元/kWh，某段原计划100 kWh，最终交付80 kWh，非紧急费用为90元。若先改为80再改为90 kWh，则按最终数量结算95元；不同版本不重复累计费用。取消电量不再进入供给，式\eqref{eq:balance}中的$q_t$改为$a_t$。

\subsubsection{预报更新与剩余计划优化}
0:00利用当时发布的光伏预报和历史负载预测，按问题二的日前框架确定原计划。6:00、12:00和18:00采用新预报时，重新构造对应的历史误差情景，以实际储电量作为初值，优化当天尚未执行的购电量。已执行区间保持不变，原计划$g_t$始终保留为结算基准。

记当天剩余区间为$\mathcal T_0$，次日前瞻区间为$\mathcal T_1$。省略已发生的固定费用，修订模型的目标为
\begin{align}
\min J_3={}&\sum_{t\in\mathcal T_0}
  \left[p_tg_t-0.5p_t\delta_t^-+1.5p_t\delta_t^+\right]
  +\sum_{t\in\mathcal T_1}p_tg_t\\
&+\sum_s\pi_s\sum_{t\in\mathcal T_0\cup\mathcal T_1}5p_te_t^{(s)}.
\end{align}
各情景共同使用当天调整量，物理约束仍采用统一模型，末端规划储电量为6000 kWh。次日购电仅用于前瞻，按普通电价计入，不预设未来必然发生新的调整。

每次求解都把维持当前计划作为候选。若预计费用改善不超过0.01元，则保持当前计划；否则在费用最优的方案中优先减少变更电量。这一比较已经包含调整费用，不再额外扣除一次违约费。

\subsubsection{滚动执行}
购电量仅在允许的发布时间修订，储能则每十分钟重新优化并执行当前动作。因此，6:10虽然可以根据新观测调整充放电，却不能再次修改6:00锁定的购电版本。控制器沿用问题二的供需规则与跨日状态递推，实际费用最终由式\eqref{eq:q3bill}计算。

这一安排使新预报先影响剩余计划，再通过储能应对区间内实际偏差。滚动优化提供的是持续修正决策的机制，并不保证每次更新后的实际费用都下降，是否有收益仍需连续回测检验。

\subsubsection{全年对照与预报时刻的价值}

对6、12、18点预报的全部八种选用组合进行独立连续回测，0点预报始终采用。另设接受全部预报但仅调整储能的策略，以区分信息更新和修改购电量的作用。各策略均保留不调整的选择，所有费用由最终交付量与实际紧急购电结算。

\begin{table}[htbp]\centering\scriptsize
\caption{问题三各策略的334天实际运行结果}\label{tab:q3compare}
\resizebox{\textwidth}{!}{\begin{tabular}{lrrrrr}
\toprule
策略 & 费用/万元 & 紧急电量/kWh & 天数 & 未利用电量/kWh & 期末储电/kWh \\
\midrule
仅0点预报 & 1468.12 & 301147.22 & 208 & 2339713.05 & 8314.11 \\
6点调整 & 1424.68 & 196735.01 & 186 & 1983450.36 & 8314.11 \\
12点调整 & 1407.42 & 192857.14 & 195 & 1871994.03 & 8359.44 \\
18点调整 & 1422.88 & 197939.78 & 208 & 2327996.23 & 8486.34 \\
6、12点调整 & 1396.07 & 166577.57 & 186 & 1738038.20 & 8359.44 \\
6、18点调整 & 1401.84 & 145568.06 & 188 & 1978655.83 & 8493.72 \\
12、18点调整 & 1394.77 & 163326.91 & 195 & 1868391.18 & 8493.72 \\
6、12、18点调整 & 1382.99 & 136287.62 & 186 & 1734278.82 & 8493.72 \\
全预报仅调储能 & 1468.55 & 307111.99 & 218 & 2353632.40 & 8314.11 \\
原计划照付口径 & 1400.86 & 127399.48 & 166 & 2439450.28 & 8493.72 \\
\bottomrule
\end{tabular}}
\end{table}
三个时刻均采用的主方案总费用为\QthreeCost\ 万元，紧急购电量为\QthreeEmergency\ kWh，发生于\QthreeDays\ 天。以仅采用0点预报为基准，费用减少\QthreeDifference\ 元，降幅为\QthreeSaving\%。费用构成为原计划金额\QthreePlanCost\ 万元，扣除取消部分退款\QthreeRefund\ 万元，加上违约费\QthreePenalty\ 万元、增购费\QthreeExtra\ 万元和紧急购电费\QthreeEmergencyCost\ 万元。

仅更新预报并调整储能，相对0点基准的实际费用增加4284.87元。进一步允许修改购电量，相对仅更新储能的方案，费用减少855533.38元。本次对照中，购电修订提供的灵活性是费用改善的主要来源。这是不同连续运行策略之间的费用差，包含其后续储电状态变化的影响。
在八种预报组合中，本次回测费用最低的是“6、12、18点调整”，总费用为1382.9943万元。该排序是样本内的运行比较，不作为重新选择主方案参数的依据。在其余预报组合固定时，加入6点预报的全年费用减少量介于113525.26至434431.22元。在其余预报组合固定时，加入12点预报的全年费用减少量介于188471.60至606987.13元。在其余预报组合固定时，加入18点预报的全年费用减少量介于126510.18至452429.31元。因此，在本次数据与模型设定下，采用6、12、18点预报并允许购电调整具有费用优势；接收预报不意味着每次都必须修改计划。在原计划照付的另一结算口径下，全部预报方案相对无调整基准的费用减少672587.38元，全年减购量为0.0000 kWh。

\begin{figure}[htbp]\centering
\includegraphics[width=\textwidth,height=.65\textheight,keepaspectratio]{figures/q3_compare.pdf}
\caption{预报选用及调整策略的费用与紧急购电比较。}
\end{figure}
\begin{figure}[htbp]\centering
\includegraphics[width=\textwidth]{figures/q3_year.pdf}
\caption{主方案逐月购电费用构成。净基础费用为原计划金额减去取消部分退款。}
\end{figure}

\subsubsection{指定日期的执行结果}
表\ref{tab:q3selected}至表\ref{tab:q3emergency}给出四个指定日期的计划、最终调整量及实际储能运行。调整后电量指非紧急交付量，加上紧急电量才是总购电量。
\begin{table}[htbp]\centering\scriptsize
\caption{指定日期的电量、总费用及首尾储电量（电量单位：kWh）}\label{tab:q3selected}
\resizebox{\textwidth}{!}{\begin{tabular}{lrrrrrr}
\toprule
日期 & 原计划电量 & 调整后电量 & 紧急电量 & 总费用/元 & 日初储电 & 日末储电 \\
\midrule
03-20 & 65824.81 & 61938.11 & 476.95 & 42591.41 & 8467.52 & 4646.88 \\
06-21 & 42586.35 & 42325.44 & 0.00 & 24691.61 & 5592.84 & 10670.88 \\
09-23 & 74771.47 & 68804.75 & 455.94 & 47093.43 & 7716.50 & 8166.23 \\
12-21 & 98420.21 & 96302.01 & 0.00 & 63767.32 & 8177.36 & 8256.55 \\
\bottomrule
\end{tabular}}
\end{table}
\begin{table}[htbp]\centering\scriptsize
\caption{指定时段的原计划与最终调整电量（kWh）}\label{tab:q3purchase}
\resizebox{\textwidth}{!}{\begin{tabular}{lrrrrrrrr}
\toprule
时段 & \multicolumn{2}{c}{03-20} & \multicolumn{2}{c}{06-21} & \multicolumn{2}{c}{09-23} & \multicolumn{2}{c}{12-21} \\
\midrule
 & 原计划 & 调整后 & 原计划 & 调整后 & 原计划 & 调整后 & 原计划 & 调整后 \\
10:00--10:10 & 0.00 & 0.00 & 0.00 & 0.00 & 0.00 & 0.00 & 429.50 & 364.56 \\
12:00--12:10 & 571.96 & 501.26 & 0.00 & 0.00 & 624.53 & 0.00 & 965.03 & 951.01 \\
14:00--14:10 & 0.00 & 0.00 & 0.00 & 0.00 & 0.00 & 0.00 & 0.00 & 0.00 \\
16:00--16:10 & 692.18 & 512.10 & 108.82 & 108.82 & 613.48 & 609.17 & 791.37 & 823.78 \\
18:00--18:10 & 667.80 & 710.59 & 543.40 & 538.05 & 744.65 & 744.65 & 710.34 & 679.21 \\
20:00--20:10 & 0.00 & 0.00 & 0.00 & 0.00 & 0.00 & 0.00 & 0.00 & 0.00 \\
\bottomrule
\end{tabular}}
\end{table}
\begin{table}[htbp]\centering\scriptsize
\caption{指定日期各四小时区间的实际充放电量（kWh）}\label{tab:q3storage}
\resizebox{\textwidth}{!}{\begin{tabular}{lrrrrrrrr}
\toprule
时段 & \multicolumn{2}{c}{03-20} & \multicolumn{2}{c}{06-21} & \multicolumn{2}{c}{09-23} & \multicolumn{2}{c}{12-21} \\
\midrule
 & 充电 & 放电 & 充电 & 放电 & 充电 & 放电 & 充电 & 放电 \\
00:00--04:00 & 1478.05 & 183.06 & 116.38 & 54.80 & 2778.41 & 109.90 & 1716.84 & 79.29 \\
04:00--08:00 & 976.35 & 5858.46 & 58.68 & 3561.88 & 881.63 & 5997.91 & 2528.56 & 1896.21 \\
08:00--12:00 & 4366.52 & 2487.31 & 3767.71 & 0.00 & 5150.47 & 1896.62 & 5245.03 & 5659.53 \\
12:00--16:00 & 6171.66 & 254.72 & 6307.99 & 0.00 & 4602.61 & 520.92 & 5590.45 & 2220.11 \\
16:00--20:00 & 247.22 & 5398.14 & 10.35 & 4681.16 & 362.32 & 5546.48 & 11.08 & 5566.55 \\
20:00--24:00 & 3769.46 & 3034.39 & 8694.40 & 2485.90 & 7769.05 & 2974.46 & 7661.63 & 2937.45 \\
\bottomrule
\end{tabular}}
\end{table}
\begin{table}[htbp]\centering\small
\caption{指定日期连续紧急购电区间及电量}\label{tab:q3emergency}
\begin{tabular}{llr}
\toprule
日期 & 紧急购电区间 & 电量/kWh \\
\midrule
03-20 & 01:20--01:30 & 12.56 \\
 & 09:10--10:00 & 290.07 \\
 & 15:40--16:40 & 169.83 \\
 & 16:50--17:00 & 4.49 \\
06-21 & 无 & 0.00 \\
09-23 & 17:10--17:20 & 24.65 \\
 & 18:20--18:40 & 38.40 \\
 & 20:30--20:40 & 28.47 \\
 & 20:50--21:40 & 364.42 \\
12-21 & 无 & 0.00 \\
\bottomrule
\end{tabular}
\end{table}


\begin{figure}[htbp]\centering
\includegraphics[width=\textwidth,height=.83\textheight,keepaspectratio]{figures/q3_day_20250320.pdf}
\caption{03-20的预报更新、购电调整与储能响应。}
\end{figure}

